\documentclass[12pt,a4paper]{article}
\usepackage{fontspec}
\setmainfont{Times New Roman}[Path=C:/Windows/Fonts/,UprightFont=times.ttf,BoldFont=timesbd.ttf,ItalicFont=timesi.ttf,BoldItalicFont=timesbi.ttf]
\usepackage{amsmath,amssymb,amsfonts}
\usepackage[left=1.5cm,right=1.5cm,top=1.5cm,bottom=1.7cm,headheight=17pt]{geometry}
\usepackage{enumitem,array,graphicx,tikz,fancyhdr,needspace,booktabs}

\pagestyle{fancy}
\fancyhf{}
\lhead{\small\textbf{ĐỀ CƯƠNG ÔN TẬP CHƯƠNG V (ĐỀ 2) -- PHƯƠNG PHÁP TỌA ĐỘ TRONG KHÔNG GIAN}}
\rhead{\small\textit{Biên soạn: Dương Gia Bảo}}
\rfoot{\small Trang \thepage}
\renewcommand{\headrulewidth}{0.4pt}
\fancypagestyle{firstpage}{\fancyhead{}\renewcommand{\headrulewidth}{0pt}}

\setlength{\parindent}{0pt}
\setlength{\parskip}{3pt}
\setlength{\emergencystretch}{1em}

\setlist[enumerate,1]{label=\textbf{Câu \arabic*:},leftmargin=*,itemsep=6pt,topsep=4pt}

% Macro trắc nghiệm 4 phương án (giãn cách 8pt giữa hai hàng theo chuẩn nhánh)
\newcommand{\bonpa}[4]{\par\nobreak\smallskip\nobreak\noindent
\begin{tabular}{@{}p{0.48\linewidth}p{0.48\linewidth}@{}}
\textbf{A.} #1 & \textbf{B.} #2\\[8pt]
\textbf{C.} #3 & \textbf{D.} #4\end{tabular}}

% Macro khung ô li viết bài làm trực tiếp (kích thước rộng rãi 5mm x 5mm)
\newcommand{\khungOli}[1][5.2cm]{%
  \needspace{2.5cm}%
  \par\nobreak\vspace{4pt}\noindent
  \begin{tikzpicture}
    \pgfmathsetmacro{\gridw}{\linewidth - 1.5pt}
    \draw[step=0.5cm, line width=0.2pt, gray!32] (0,0) grid (\gridw pt, #1);
    \draw[line width=0.6pt, gray!60] (0,0) rectangle (\gridw pt, #1);
    \node[anchor=north east, font=\footnotesize\itshape\color{gray!60}, inner sep=4pt] at (\gridw pt, #1) {Bài làm của học sinh};
  \end{tikzpicture}
  \par\vspace{7pt}
}

\begin{document}
\thispagestyle{firstpage}

\noindent
\begin{minipage}[t]{0.36\textwidth}
    \centering
    {\fontsize{12}{14}\selectfont\textbf{ĐỀ CƯƠNG ÔN TẬP}}\\[4pt]
    \textit{Biên soạn: Dương Gia Bảo}\\[6pt]
    \setlength{\fboxrule}{0.8pt}
    \setlength{\fboxsep}{5pt}
    \fbox{\begin{minipage}{0.88\linewidth}
        \centering
        {\fontsize{12}{14}\selectfont\textbf{CHƯƠNG V -- ĐỀ 2}}
    \end{minipage}}
\end{minipage}\hfill
\begin{minipage}[t]{0.62\textwidth}
    \centering
    {\fontsize{13}{15}\selectfont\textbf{ĐỀ CƯƠNG ÔN TẬP CHƯƠNG V -- ĐỀ SỐ 2}}\\[3pt]
    {\fontsize{12}{14}\selectfont\textbf{PHƯƠNG PHÁP TỌA ĐỘ TRONG KHÔNG GIAN}}\\[3pt]
    {\fontsize{11}{13}\selectfont\textit{(GDPT 2018 -- 20 Trắc nghiệm, 5 Đúng/Sai, 25 Trả lời ngắn)}}
\end{minipage}

\par\vspace{4pt}
\hrule height 0.6pt
\vspace{4pt}
\noindent\textit{Họ và tên học sinh:}\dotfill\hspace{0.8cm}\textit{Lớp:}\makebox[3.2cm]{\dotfill}
\par\vspace{3pt}
\hrule height 0.6pt
\vspace{6pt}

%% ==============================================================
%% PHẦN I: TRẮC NGHIỆM NHIỀU LỰA CHỌN (20 CÂU)
%% ==============================================================
\begin{center}
    {\fontsize{13}{15}\selectfont\textbf{PHẦN I. CÂU TRẮC NGHIỆM NHIỀU LỰA CHỌN (20 CÂU)}}
\end{center}
\textit{Thí sinh trả lời từ Câu 1 đến Câu 20. Mỗi câu hỏi chỉ chọn một phương án.}

\begin{enumerate}
\item Trong không gian $Oxyz$, điểm nào dưới đây thuộc mặt phẳng $(P): 3x - 2y + z - 4 = 0$?
\bonpa{$A(1; 0; 1)$}{$B(1; 1; 1)$}{$C(2; 1; 1)$}{$D(0; 2; 1)$}

\item Trong không gian $Oxyz$, phương trình mặt phẳng đi qua gốc tọa độ $O(0; 0; 0)$ và song song với mặt phẳng $(\alpha): 2x - y + 5z - 7 = 0$ là
\bonpa{$2x - y + 5z - 7 = 0$}{$2x - y + 5z = 0$}{$2x + y + 5z = 0$}{$2x - y - 5z = 0$}

\item Trong không gian $Oxyz$, khoảng cách từ điểm $A(2; -1; 3)$ đến mặt phẳng $(P): x + 2y - 2z + 12 = 0$ bằng
\bonpa{$1$}{$3$}{$2$}{$4$}

\item Trong không gian $Oxyz$, mặt phẳng đi qua ba điểm $M(3; 0; 0)$, $N(0; 2; 0)$, $P(0; 0; -1)$ có phương trình tổng quát là
\bonpa{$2x + 3y + 6z - 6 = 0$}{$2x - 3y - 6z - 6 = 0$}{$3x + 2y - z - 6 = 0$}{$2x + 3y - 6z - 6 = 0$}

\item Trong không gian $Oxyz$, cho hai mặt phẳng $(P): 2x - 3y + z - 1 = 0$ và $(Q): 4x - 6y + 2z - 2 = 0$. Vị trí tương đối của hai mặt phẳng $(P)$ và $(Q)$ là
\bonpa{$(P)$ trùng với $(Q)$}{$(P)$ song song với $(Q)$}{$(P)$ cắt và vuông góc với $(Q)$}{$(P)$ cắt nhưng không vuông góc với $(Q)$}

\item Trong không gian $Oxyz$, đường thẳng đi qua hai điểm $A(1; 2; -1)$ và $B(2; 4; 1)$ có một vectơ chỉ phương là
\bonpa{$\vec{u}_1 = (3; 6; 0)$}{$\vec{u}_2 = (1; 2; 2)$}{$\vec{u}_3 = (-1; -2; 2)$}{$\vec{u}_4 = (1; 2; 0)$}

\item Trong không gian $Oxyz$, phương trình tham số của đường thẳng $d$ đi qua điểm $M(2; -3; 1)$ và có vectơ chỉ phương $\vec{u} = (-1; 2; 3)$ là
\bonpa{$\begin{cases} x = -1 + 2t \\ y = 2 - 3t \\ z = 3 + t \end{cases}$}{$\begin{cases} x = 2 - t \\ y = 3 + 2t \\ z = 1 + 3t \end{cases}$}{$\begin{cases} x = 2 - t \\ y = -3 + 2t \\ z = 1 + 3t \end{cases}$}{$\begin{cases} x = 2 + t \\ y = -3 - 2t \\ z = 1 - 3t \end{cases}$}

\item Trong không gian $Oxyz$, cho đường thẳng $d: \dfrac{x+1}{2} = \dfrac{y-3}{-2} = \dfrac{z}{1}$ và mặt phẳng $(P): 2x - 2y + z + 3 = 0$. Khẳng định nào sau đây đúng?
\bonpa{$d$ song song với $(P)$}{$d$ nằm trong $(P)$}{$d$ cắt nhưng không vuông góc với $(P)$}{$d$ vuông góc với $(P)$}

\item Trong không gian $Oxyz$, hình chiếu vuông góc của điểm $A(3; -4; 5)$ lên trục tọa độ $Oz$ có tọa độ là
\bonpa{$(0; 0; 5)$}{$(3; -4; 0)$}{$(3; 0; 0)$}{$(0; -4; 0)$}

\item Trong không gian $Oxyz$, cho đường thẳng $d: \dfrac{x-1}{2} = \dfrac{y+1}{1} = \dfrac{z-2}{3}$. Điểm nào sau đây KHÔNG thuộc đường thẳng $d$?
\bonpa{$M(1; -1; 2)$}{$N(2; 1; 3)$}{$P(3; 0; 5)$}{$Q(-1; -2; -1)$}

\item Trong không gian $Oxyz$, cho hai đường thẳng $d_1: \dfrac{x-1}{1} = \dfrac{y}{1} = \dfrac{z+2}{0}$ và $d_2: \dfrac{x+1}{0} = \dfrac{y-2}{1} = \dfrac{z}{1}$. Góc giữa hai đường thẳng $d_1$ và $d_2$ bằng
\bonpa{$30^\circ$}{$45^\circ$}{$60^\circ$}{$90^\circ$}

\item Trong không gian $Oxyz$, góc giữa hai mặt phẳng tọa độ $(Oxy)$ và $(Oyz)$ bằng
\bonpa{$0^\circ$}{$30^\circ$}{$45^\circ$}{$90^\circ$}

\item Trong không gian $Oxyz$, góc giữa đường thẳng $d: \dfrac{x}{1} = \dfrac{y}{\sqrt{3}} = \dfrac{z}{0}$ và trục tọa độ $Ox$ bằng
\bonpa{$60^\circ$}{$30^\circ$}{$45^\circ$}{$90^\circ$}

\item Trong không gian $Oxyz$, sin của góc giữa đường thẳng $d: \dfrac{x-2}{2} = \dfrac{y+1}{1} = \dfrac{z-1}{2}$ và mặt phẳng $(P): x + y + z - 5 = 0$ bằng
\bonpa{$\dfrac{5}{3}$}{$\dfrac{5\sqrt{3}}{9}$}{$\dfrac{\sqrt{3}}{9}$}{$\dfrac{5}{9}$}

\item Trong không gian $Oxyz$, cho hai đường thẳng $d_1: \dfrac{x}{1} = \dfrac{y-1}{m} = \dfrac{z}{2}$ và $d_2: \dfrac{x+1}{2} = \dfrac{y}{1} = \dfrac{z-3}{-1}$. Hai đường thẳng $d_1$ và $d_2$ vuông góc với nhau khi giá trị của tham số $m$ bằng
\bonpa{$m = 1$}{$m = 2$}{$m = 0$}{$m = -2$}

\item Trong không gian $Oxyz$, phương trình nào sau đây là phương trình của một mặt cầu?
\bonpa{$x^2 + y^2 + z^2 - 2x + 2y - 4z + 10 = 0$}{$x^2 + y^2 - z^2 - 2x + 4y - 6z + 1 = 0$}{$x^2 + y^2 + z^2 - 2x + 4y - 6z + 20 = 0$}{$x^2 + y^2 + z^2 - 2x + 4y - 6z + 1 = 0$}

\item Trong không gian $Oxyz$, phương trình mặt cầu tâm $I(1; -3; 2)$ và đi qua gốc tọa độ $O(0; 0; 0)$ là
\bonpa{$(x-1)^2 + (y+3)^2 + (z-2)^2 = 14$}{$(x+1)^2 + (y-3)^2 + (z+2)^2 = 14$}{$(x-1)^2 + (y+3)^2 + (z-2)^2 = \sqrt{14}$}{$(x-1)^2 + (y+3)^2 + (z-2)^2 = 12$}

\item Trong không gian $Oxyz$, cho mặt cầu $(S): (x+1)^2 + (y-2)^2 + (z+3)^2 = 25$. Mặt phẳng tiếp xúc với mặt cầu $(S)$ tại điểm $A(2; 2; 1)$ có một vectơ pháp tuyến là
\bonpa{$\vec{n}_1 = (1; 0; 4)$}{$\vec{n}_2 = (3; 0; 4)$}{$\vec{n}_3 = (3; 4; 4)$}{$\vec{n}_4 = (3; 0; -2)$}

\item Trong không gian $Oxyz$, tọa độ tâm $I$ của mặt cầu $(S): x^2 + y^2 + z^2 - 4x + 6y - 2z + 5 = 0$ là
\bonpa{$I(-2; 3; -1)$}{$I(4; -6; 2)$}{$I(2; -3; 1)$}{$I(-4; 6; -2)$}

\item Trong không gian $Oxyz$, một khinh khí cầu quan sát khí quyển đang bay ở vị trí có tọa độ $M(3; 4; 5)$ (đơn vị: km) so với trạm khí tượng mặt đất đặt tại gốc tọa độ $O(0; 0; 0)$. Khoảng cách đường thẳng từ trạm khí tượng $O$ đến khinh khí cầu $M$ bằng
\bonpa{$12\text{ km}$}{$5\text{ km}$}{$10\text{ km}$}{$5\sqrt{2}\text{ km}$}
\end{enumerate}

\newpage
%% ==============================================================
%% PHẦN II: TRẮC NGHIỆM ĐÚNG SAI (5 CÂU)
%% ==============================================================
\begin{center}
    {\fontsize{13}{15}\selectfont\textbf{PHẦN II. CÂU TRẮC NGHIỆM ĐÚNG SAI (5 CÂU)}}
\end{center}
\textit{Thí sinh trả lời từ Câu 1 đến Câu 5. Trong mỗi ý a), b), c), d) ở mỗi câu, thí sinh chọn đúng hoặc sai.}

\setlist[enumerate,1]{label=\textbf{Câu \arabic*:},leftmargin=*,itemsep=10pt,topsep=6pt}

\begin{enumerate}
\needspace{6.5cm}
\item Trong không gian $Oxyz$, cho hình hộp chữ nhật $ABCD.A'B'C'D'$ có đỉnh $A$ trùng với gốc tọa độ $O(0; 0; 0)$, đỉnh $B$ thuộc tia $Ox$, đỉnh $D$ thuộc tia $Oy$ và đỉnh $A'$ thuộc tia $Oz$. Biết tọa độ các đỉnh $B(4; 0; 0)$, $D(0; 6; 0)$ và $A'(0; 0; 3)$.
\begin{enumerate}[label=\alph*),leftmargin=*,itemsep=3pt]
    \item Tọa độ đỉnh $C$ là $(4; 6; 0)$ và đỉnh $C'$ là $(4; 6; 3)$.
    \item Phương trình tổng quát của mặt phẳng $(A'BD)$ là $3x + 2y + 4z - 12 = 0$.
    \item Thể tích của khối hộp chữ nhật $ABCD.A'B'C'D'$ bằng $72$.
    \item Khoảng cách từ điểm $C$ đến mặt phẳng $(A'BD)$ gấp đôi khoảng cách từ điểm $A$ đến mặt phẳng $(A'BD)$.
\end{enumerate}

\needspace{6.5cm}
\item Trong không gian $Oxyz$, cho đường thẳng $d: \begin{cases} x = 1 + 2t \\ y = -2 + t \\ z = 3 - t \end{cases}$ ($t \in \mathbb{R}$) và mặt phẳng $(P): x + y + 3z - 8 = 0$.
\begin{enumerate}[label=\alph*),leftmargin=*,itemsep=3pt]
    \item Điểm $M(1; -2; 3)$ thuộc đường thẳng $d$.
    \item Đường thẳng $d$ song song với mặt phẳng $(P)$.
    \item Đường thẳng $d$ nằm hoàn toàn trong mặt phẳng $(P)$.
    \item Mặt phẳng $(Q)$ đi qua điểm $A(0; 1; 2)$ và vuông góc với đường thẳng $d$ có phương trình là $2x + y - z + 1 = 0$.
\end{enumerate}

\needspace{6.5cm}
\item Trong không gian $Oxyz$, cho hai điểm $A(1; 0; 2)$, $B(3; 2; 0)$ và mặt phẳng $(\alpha): x + 2y - 2z + 6 = 0$.
\begin{enumerate}[label=\alph*),leftmargin=*,itemsep=3pt]
    \item Mặt cầu $(S)$ nhận đoạn thẳng $AB$ làm đường kính có tâm là điểm $I(2; 1; 1)$.
    \item Bán kính của mặt cầu $(S)$ bằng $\sqrt{3}$.
    \item Khoảng cách từ tâm $I$ của mặt cầu $(S)$ đến mặt phẳng $(\alpha)$ bằng $\dfrac{8}{3}$.
    \item Mặt phẳng $(\alpha)$ cắt mặt cầu $(S)$ theo giao tuyến là một đường tròn.
\end{enumerate}

\needspace{6.5cm}
\item Trong không gian $Oxyz$, cho hai đường thẳng $d_1: \dfrac{x-1}{1} = \dfrac{y+1}{-2} = \dfrac{z-2}{2}$ và $d_2: \dfrac{x+2}{2} = \dfrac{y-3}{1} = \dfrac{z}{-1}$.
\begin{enumerate}[label=\alph*),leftmargin=*,itemsep=3pt]
    \item Vectơ $\vec{u}_1 = (1; -2; 2)$ là một vectơ chỉ phương của $d_1$ và $\vec{u}_2 = (2; 1; -1)$ là một vectơ chỉ phương của $d_2$.
    \item Hai đường thẳng $d_1$ và $d_2$ vuông góc với nhau.
    \item Côsin của góc giữa hai đường thẳng $d_1$ và $d_2$ bằng $\dfrac{\sqrt{6}}{9}$.
    \item Hai đường thẳng $d_1$ và $d_2$ cắt nhau tại một điểm.
\end{enumerate}

\needspace{6.5cm}
\item Trong hệ trục toạ độ $Oxyz$ mô phỏng không gian vũ trụ quanh Trái Đất (gốc $O$ đặt tại tâm Trái Đất, bán kính Trái Đất $R_0 = 6400\text{ km}$, đơn vị đo trên các trục tọa độ là $1000\text{ km}$). Một vệ tinh nhân tạo địa tĩnh đang ở vị trí cố định $S(0; 0; 42)$ và một trạm quan sát khí tượng đặt tại điểm $A(0; 3{,}84; 5{,}12)$.
\begin{enumerate}[label=\alph*),leftmargin=*,itemsep=3pt]
    \item Trạm quan sát khí tượng $A$ nằm trên bề mặt khối cầu Trái Đất.
    \item Khoảng cách từ tâm Trái Đất $O$ đến vệ tinh $S$ bằng $42000\text{ km}$.
    \item Khoảng cách truyền tín hiệu đường thẳng từ trạm quan sát $A$ đến vệ tinh $S$ nhỏ hơn $38000\text{ km}$.
    \item Góc nâng quan sát từ trạm $A$ đến vệ tinh $S$ (góc giữa tia nhìn $AS$ và mặt phẳng tiếp diện với Trái Đất tại $A$) lớn hơn $80^\circ$.
\end{enumerate}
\end{enumerate}

\newpage
%% ==============================================================
%% PHẦN III: TRẢ LỜI NGẮN (25 CÂU)
%% ==============================================================
\begin{center}
    {\fontsize{13}{15}\selectfont\textbf{PHẦN III. CÂU TRẮC NGHIỆM TRẢ LỜI NGẮN KÈM Ô LI (25 CÂU)}}
\end{center}
\textit{Thí sinh trình bày lời giải và điền đáp số vào khung ô li làm bài tương ứng.}

\setlist[enumerate,1]{label=\textbf{Câu \arabic*:},leftmargin=*,itemsep=6pt,topsep=4pt}

\subsubsection*{Mạch 1. Phương trình mặt phẳng và khoảng cách (Câu 1 -- 6)}

\begin{enumerate}
\needspace{7.6cm}
\item Trong không gian $Oxyz$, cho ba điểm $A(1; 0; 0)$, $B(0; 2; 0)$, $C(0; 0; 3)$. Mặt phẳng $(P)$ song song với mặt phẳng $(ABC)$ và đi qua điểm $M(1; 1; 1)$ có phương trình dạng $6x + 3y + 2z + D = 0$. Tìm giá trị của hằng số $D$.
\khungOli[5.2cm]

\needspace{7.6cm}
\item Trong không gian $Oxyz$, tính khoảng cách từ điểm $A(1; 2; 3)$ đến mặt phẳng $(P): 3x - 4z + 14 = 0$.
\khungOli[5.2cm]

\needspace{7.6cm}
\item Trong không gian $Oxyz$, mặt phẳng $(P)$ đi qua hai điểm $A(1; 1; 0)$, $B(0; 2; 1)$ và vuông góc với mặt phẳng $(Q): x + y - z + 2 = 0$. Biết phương trình mặt phẳng $(P)$ có dạng $Ax + By + z + D = 0$. Tính giá trị của biểu thức $S = A + B + D$.
\khungOli[5.2cm]

\needspace{7.6cm}
\item Trong không gian $Oxyz$, tính khoảng cách giữa hai mặt phẳng song song $(P): x - 2y + 2z - 1 = 0$ và $(Q): x - 2y + 2z + 8 = 0$.
\khungOli[5.2cm]

\needspace{7.6cm}
\item Trong không gian $Oxyz$, tìm giá trị dương của tham số $m$ để khoảng cách từ gốc tọa độ $O(0; 0; 0)$ đến mặt phẳng $(P): 2x - y + 2z + m = 0$ bằng $4$.
\khungOli[5.2cm]

\needspace{7.8cm}
\item Trong không gian $Oxyz$, mặt phẳng trần của một căn phòng có phương trình $2x + y + 2z - 18 = 0$, sàn nhà nằm trên mặt phẳng tọa độ $(Oxy): z = 0$ (đơn vị: mét). Một chiếc quạt trần được treo tại điểm $I(1; 2; 4)$. Tính khoảng cách từ chiếc quạt trần $I$ đến mặt phẳng sàn nhà (đơn vị: mét).
\khungOli[5.2cm]

\subsubsection*{Mạch 2. Phương trình đường thẳng và vị trí tương đối (Câu 7 -- 12)}

\needspace{7.6cm}
\item Trong không gian $Oxyz$, tìm tọa độ giao điểm $M(x_0; y_0; z_0)$ của đường thẳng $d: \dfrac{x-3}{1} = \dfrac{y+1}{-2} = \dfrac{z-2}{1}$ và mặt phẳng $(P): x + 2y + z + 1 = 0$. Tính giá trị của biểu thức $T = x_0 + y_0 + z_0$.
\khungOli[5.2cm]

\needspace{7.6cm}
\item Trong không gian $Oxyz$, cho điểm $A(2; 1; 3)$ và đường thẳng $d: \dfrac{x-1}{1} = \dfrac{y-2}{1} = \dfrac{z-3}{1}$. Gọi $H(x_0; y_0; z_0)$ là hình chiếu vuông góc của điểm $A$ lên đường thẳng $d$. Tính giá trị của $S = x_0 + y_0 + z_0$.
\khungOli[5.2cm]

\needspace{7.6cm}
\item Trong không gian $Oxyz$, cho điểm $A(2; 1; -3)$. Gọi $A'$ là điểm đối xứng với điểm $A$ qua mặt phẳng tọa độ $(Oxy)$. Tính độ dài đoạn thẳng $AA'$.
\khungOli[5.2cm]

\needspace{7.8cm}
\item Trong không gian $Oxyz$, cho hai đường thẳng $d_1: \dfrac{x-1}{2} = \dfrac{y-2}{3} = \dfrac{z-3}{4}$ và $d_2: \dfrac{x-3}{4} = \dfrac{y-5}{6} = \dfrac{z-7}{8}$. Hãy xác định vị trí tương đối của hai đường thẳng này. (Quy ước: Nếu $d_1$ trùng $d_2$, ghi kết quả bằng $0$. Nếu $d_1$ song song $d_2$, tính khoảng cách giữa hai đường thẳng).
\khungOli[5.2cm]

\needspace{7.6cm}
\item Trong không gian $Oxyz$, cho tam giác $ABC$ có tọa độ ba đỉnh lần lượt là $A(1; 2; 3)$, $B(2; 1; 0)$ và $C(3; 3; 3)$. Gọi $G(x_G; y_G; z_G)$ là tọa độ trọng tâm của tam giác $ABC$. Tính tổng hoành độ, tung độ và cao độ $x_G + y_G + z_G$.
\khungOli[5.2cm]

\needspace{8.5cm}
\item Trong không gian $Oxyz$, cho hai đường thẳng chéo nhau:
\[
d_1: \begin{cases} x = 1 + t \\ y = 0 \\ z = 2 \end{cases}\quad\text{và}\quad d_2: \begin{cases} x = 0 \\ y = t' \\ z = -1. \end{cases}
\]
Tính khoảng cách ngắn nhất giữa hai đường thẳng $d_1$ và $d_2$.
\khungOli[5.2cm]

\subsubsection*{Mạch 3. Công thức tính góc trong không gian (Câu 13 -- 17)}

\needspace{7.6cm}
\item Trong không gian $Oxyz$, tính số đo góc (theo độ, $^\circ$) giữa đường thẳng $d: \dfrac{x-1}{1} = \dfrac{y+2}{1} = \dfrac{z}{0}$ và trục tọa độ $Oy$.
\khungOli[5.2cm]

\needspace{7.6cm}
\item Trong không gian $Oxyz$, tính số đo góc (theo độ, $^\circ$) giữa đường thẳng $d: \dfrac{x}{\sqrt{3}} = \dfrac{y}{1} = \dfrac{z}{0}$ và mặt phẳng $(P): z + 5 = 0$.
\khungOli[5.2cm]

\needspace{7.6cm}
\item Trong không gian $Oxyz$, côsin của góc giữa hai mặt phẳng $(P): x + y - 4 = 0$ và $(Q): y + z + 2 = 0$ viết dưới dạng phân số tối giản $\dfrac{a}{b}$ với $a, b \in \mathbb{Z}^+$. Tính giá trị của $a + b$.
\khungOli[5.2cm]

\needspace{7.6cm}
\item Trong không gian $Oxyz$, cho đường thẳng $d_1$ có vectơ chỉ phương $\vec{u}_1 = (1; 0; m)$ và đường thẳng $d_2$ có vectơ chỉ phương $\vec{u}_2 = (1; 1; 0)$. Tìm giá trị dương của tham số $m$ để góc giữa hai đường thẳng $d_1$ và $d_2$ bằng $60^\circ$.
\khungOli[5.2cm]

\needspace{7.8cm}
\item Một tấm pin năng lượng mặt trời đặt trên mái nhà nghiêng có phương trình mặt phẳng là $x + y + \sqrt{2}z - 10 = 0$. Mặt đất phẳng nằm trên mặt phẳng tọa độ $(Oxy): z = 0$. Tính góc nghiêng (theo độ, $^\circ$) giữa tấm pin năng lượng mặt trời và mặt phẳng nằm ngang.
\khungOli[5.2cm]

\subsubsection*{Mạch 4. Phương trình mặt cầu và vị trí tương đối (Câu 18 -- 22)}

\needspace{7.6cm}
\item Trong không gian $Oxyz$, tính bán kính $R$ của mặt cầu $(S): x^2 + y^2 + z^2 - 4x + 2y - 20 = 0$.
\khungOli[5.2cm]

\needspace{7.6cm}
\item Trong không gian $Oxyz$, cho mặt cầu $(S)$ có tâm $I(2; 1; -1)$ và tiếp xúc với mặt phẳng $(P): 2x - 2y + z + 8 = 0$. Gọi $S$ là diện tích toàn phần của mặt cầu $(S)$. Hãy ghi giá trị của biểu thức $\dfrac{S}{\pi}$.
\khungOli[5.2cm]

\needspace{7.6cm}
\item Trong không gian $Oxyz$, mặt phẳng $(P): z - 2 = 0$ cắt mặt cầu $(S): x^2 + y^2 + z^2 = 13$ theo giao tuyến là một đường tròn có chu vi $C$. Hãy ghi giá trị của biểu thức $\dfrac{C}{\pi}$.
\khungOli[5.2cm]

\needspace{7.8cm}
\item Trong một nhà xưởng sản xuất, một trạm phát sóng WiFi nội bộ đặt tại điểm $I(2; 3; 2)$ có bán kính vùng phủ sóng ổn định bằng $6\text{ m}$. Một robot tự hành chuyển động thẳng đều dọc theo đường ray dẫn hướng $d: \begin{cases} x = 2 \\ y = -5 + t \\ z = 2 \end{cases}$ ($t \in \mathbb{R}$, đơn vị: mét). Tính quãng đường (đơn vị: mét) mà robot di chuyển trong phạm vi nhận được tín hiệu sóng WiFi ổn định từ trạm phát.
\khungOli[5.2cm]

\needspace{7.6cm}
\item Trong không gian $Oxyz$, cho mặt cầu $(S): (x-1)^2 + (y+1)^2 + z^2 = 4$. Tìm giá trị dương của tham số $m$ để mặt phẳng $(P): 2x - y + 2z + m = 0$ tiếp xúc với mặt cầu $(S)$.
\khungOli[5.2cm]

\subsubsection*{Mạch 5. Bài toán tổng hợp và cực trị hình học Oxyz (Câu 23 -- 25)}

\needspace{7.6cm}
\item Trong không gian $Oxyz$, cho hai điểm $A(1; 1; 1)$ và $B(3; 3; 3)$. Điểm $M(x_M; 0; 0)$ nằm trên trục tọa độ $Ox$ sao cho tổng bình phương khoảng cách $MA^2 + MB^2$ đạt giá trị nhỏ nhất. Tìm hoành độ $x_M$ của điểm $M$.
\khungOli[5.2cm]

\needspace{7.6cm}
\item Trong không gian $Oxyz$, cho điểm $A(3; 2; 2)$ và mặt phẳng $(P): x + y + z - 1 = 0$. Gọi $A'$ là điểm đối xứng với điểm $A$ qua mặt phẳng $(P)$. Tính khoảng cách từ điểm $A'$ đến gốc tọa độ $O(0; 0; 0)$.
\khungOli[5.2cm]

\needspace{7.8cm}
\item Trong một công trình xây dựng, đầu vươn của một cánh tay cẩu có thể chuyển động chạm tới các điểm $M(x; y; z)$ nằm trên bề mặt khối cầu an toàn $(S): x^2 + y^2 + (z - 10)^2 = 25$ (đơn vị trên các trục tọa độ là mét). Mặt đất phẳng được quy ước là mặt phẳng $(Oxy): z = 0$. Tìm độ cao lớn nhất $z_{\max}$ (đơn vị: mét) so với mặt đất mà đầu cần cẩu có thể vươn tới.
\khungOli[5.2cm]
\end{enumerate}

\end{document}
